\section*{Introduction}

Language models are rarely used bare. Deployed assistants
run behind a system prompt that gives them a role, such
as a tutor, a customer support agent, or a medical or
security assistant. A role is meant to change tone and
scope, not the model's safety behavior. Safety behavior
is, however, learned from human feedback
\citep{ouyang2022training, bai2022training}, and a model
conditioned on a role may infer properties of the agent
that would produce such text and generate accordingly
\citep{andreas-2022-language}. Whether a role leaves
refusal behavior untouched is therefore an empirical
question.

Refusal can fail in two ways. A model that refuses too
little complies with harmful requests. A model that
refuses too much turns away benign requests that merely
sound dangerous, such as ``How do I kill a Python
process?'' \citet{rottger-etal-2024-xstest} introduced
XSTest to measure this second failure, which they call
exaggerated safety, using safe prompts alongside unsafe
contrast prompts. A model is well calibrated in this
sense if it complies with the former and refuses the
latter.\footnote{We use calibration for refusal
\emph{behavior}; it is unrelated to the calibration
of confidence scores.}

Persona prompts have mostly been studied from two
angles. On the attack side, optimized persona prompts
can lower refusal rates by 50--70\% across several
models \citep{zhang2025enhancing}, while expert personas
have been reported to improve safety alignment at some
cost in accuracy \citep{hu2026expert}. On the bias side,
\citet{plaza-del-arco-etal-2025-no} varied 15
sociodemographic personas across 16 models and three
classification tasks and found that persona effects
on false refusal depend on the model and the task,
shrink for more capable models, and may have been
overestimated. What is missing is a controlled test
of the kind of persona used in deployed assistants,
an occupational or role persona, on both sides of
refusal calibration at once. A persona that makes a
model refuse benign requests and one that makes it
comply with unsafe ones call for different remedies,
and a single refusal rate cannot tell them apart.

A second difficulty is measurement. Responses are free
text, so comparing conditions requires labeling them
as compliance or refusal, which is often delegated
to a language model judge. Such labels can be wrong
in ways that depend on response style, and style is
exactly what a persona changes. Evaluating persona
effects therefore also requires checking the
reliability of these judgments.

We ask: \emph{Does persona prompting change an
instruction-tuned model's over-refusal and under-refusal
on XSTest, and does it change the model's ability to
tell safe from unsafe requests?} We investigate two
hypotheses. \textbf{H1}: personas differ in their rates
of over-refusal on safe prompts and of under-refusal
on unsafe prompts. \textbf{H2}: personas differ in how
well they separate safe from unsafe prompts, not
merely in their overall refusal level. Following the
logic of behavioral assessment of language models
\citep[e.g.,][]{binz2023using}, we evaluate
Qwen2.5-7B-Instruct
\citep{qwen2025qwen25technicalreport} under a neutral
baseline and five role personas, using a systematic
benchmark rather than anecdotes.

Our contributions are: (i) a controlled comparison
of role personas on both over-refusal and under-refusal;
(ii) a labeling scheme that distinguishes off-topic
deflection from refusal, together with an audit of
LLM judge reliability; and (iii) code, data, and
annotations to support reproduction and further
analysis, available at \href{https://github.com/gursul/persona-safety-calibration}{GitHub}.
